\documentclass[10pt,a4paper]{amsart}
\usepackage[margin=19mm]{geometry}
\usepackage{amsmath,amssymb,mathtools}
\usepackage[scr=rsfs,cal=euler]{mathalfa}
\usepackage{xfrac}
\usepackage[unicode,hidelinks,bookmarks=false]{hyperref}
\newcommand{\ie}{i.e.\ }
\newcommand{\cf}{cf.\ }
\newcommand{\resp}{respectively\ }
\newcommand{\Subsection}{\subsection}
\providecommand{\nobreakdash}{\nobreak\hbox{-}}
\setlength{\emergencystretch}{2em}
\multlinegap0pt
\usepackage{amssymb}
\usepackage{xcolor}
\usepackage{algorithm}\usepackage{algpseudocode}
\usepackage{bm}
\usepackage{mathtools}
\newtheorem{lemma}{Lemma}
\title{On solving nonlinear simultaneous equations arising from the double-exponential Sinc-collocation method for initial value problems}
\author{Yusaku Yamamoto, Ken'ichiro Tanaka}
\date{Source: arXiv:2601.01007v2 (CC BY 4.0)}
\begin{document}
\maketitle

\noindent\emph{Source and licence.} On solving nonlinear simultaneous equations arising from the double-exponential Sinc-collocation method for initial value problems. Version arXiv:2601.01007v2 (CC BY 4.0), distributed by arXiv under the Creative Commons Attribution 4.0 International licence, http://creativecommons.org/licenses/by/4.0/, as recorded in the arXiv licence metadata for this exact version and shown on the abstract page https://arxiv.org/abs/2601.01007v2. Full licence notice: \texttt{licenses/arxiv-2601.01007-CC-BY-4.0.txt}.\par\smallskip

\noindent\textit{Editorial scope.} Counted unit: the article's lemma Lem\_diff\_DE\_monotone (source0.tex lines 1063--1071). The article's complete original proof of it is delivered with it (source0.tex lines 1073--1148, one \texttt{proof} environment of 76 lines). No mathematics is written, repaired or re-derived: the statement and the proof are the article's own lines, in the article's own notation, and no retained line is substituted or re-flowed.  No further source-local context is needed: every cross-reference of every retained line resolves inside the two delivered spans.  Nothing was omitted from any retained span, so the package carries no omission record.  No retained line contains a citation, so the wrapper carries no bibliography and no bibliography entry.  No retained line contains a figure environment, an image-inclusion command or drawing code, so no image is delivered and the package declares no asset dependency.  Editorial formatting only, in addition: an AMS \texttt{amsart} preamble that restates the article's own declarations and macros used by the retained lines, namely the environments \texttt{lemma} on separate counters, together with the standard packages those lines need.  The retained environments print in this document as Lemma 1 in order of appearance, whereas the article prints the same items with the numbering of its own full document.  The complete native source of the article as distributed in the arXiv e-print for this exact version is bundled unchanged as \texttt{sources/192041/source0.tex}.
\begin{lemma}
\label{Lem_diff_DE_monotone}
The function 
\[
f(x) =
\frac{\frac{\pi}{2}\cosh(x)}{\cosh^2\left(\frac{\pi}{2}\sinh(x)\right)}
\]
is even on $\mathbb{R}$ and monotone decreasing on $\{ x \in \mathbb{R} \mid  x \geq 0\}$. 
\end{lemma}

\begin{proof}
It is trivial that $f$ is even. 
Below we prove the monotonicity for $x \geq 0$. 
By letting $u(x) = \frac{\pi}{2} \sinh(x)$, we have
\[
f'(x) = 
\frac{\pi}{2} \cdot
\frac{\sinh(x) \cosh(u(x)) - \pi \cosh^{2}(x) \sinh(u(x))}{\cosh^{3}(u(x))}.
\]
Since the sign of the right-hand side is equal to that of its numerator, 
the inequality $f'(x) \leq 0$ is equivalent to that 
\begin{align}
\frac{\sinh(x)}{\pi \cosh^{2}(x)} \leq \tanh(u(x)). 
\notag
\end{align}
By letting $t = \sinh(x)$,
this inequality is equivalent to that
\begin{align}
\frac{t}{\pi (1+t^{2})} \leq \tanh \left( \frac{\pi}{2} t \right). 
\label{eq:ineq_equiv_concl}
\end{align}
Therefore it suffices to prove this inequality for any $t \geq 0$.

Let $g(t)$ be defined by $g(t) = \frac{t}{\pi (1+t^{2})}$ for $t \geq 0$. 
Then we have
\begin{align}
g'(t) = \frac{1-t^{2}}{\pi(1+t^{2})^{2}}
\quad \text{and} \quad
g''(t) = \frac{2t(-t^{3} + 2t^{2} - t -2)}{\pi(1+t^{2})^{3}}.
\notag
\end{align}
Therefore $g(t)$ attains its maximum value $\frac{1}{2\pi}$ at $t = 1$. 
In addition, 
for $t \in [0, 1]$, 
we have $-t^{3} + 2t^{2} - t -2 \leq -t^{3} + 2 - t - 2 = -t^{3} - t \leq 0$, 
which implies $g''(t) \leq 0$. 
Therefore $g(t)$ is concave on $[0,1]$ and we have
\begin{align}
g(t) \leq \frac{1}{\pi} t \quad (=g'(0)t).
\label{eq:key_ineq_g_t_on_0_1}
\end{align}
for $t \in [0,1]$. 
Next we show 
\begin{align}
\frac{1}{\pi} < \frac{\mathrm{e}^{\pi} - 1}{\mathrm{e}^{\pi} + 1}
\quad 
\left( = \tanh \left( \frac{\pi}{2} \right) \right).
\label{eq:key_ineq_1_pi_tanh}
\end{align}
Indeed, this inequality is equivalent to that 
\(
\mathrm{e}^{\pi} > \frac{\pi + 1}{\pi-1}
\), which is shown as
\[
\mathrm{e}^{\pi}
\geq 1 + 3 + \frac{3^2}{2!} + \frac{3^3}{3!}
= 13 > 2 = 
\frac{3+1}{3-1}
\geq 
\frac{\pi + 1}{\pi-1}. 
\]
Then, 
by combining~\eqref{eq:key_ineq_g_t_on_0_1}, \eqref{eq:key_ineq_1_pi_tanh}, 
and the concavity of $t \mapsto \tanh(\frac{\pi}{2}t)$ on $[0,1]$, 
we have
\begin{align}
g(t) \leq \frac{1}{\pi} t \leq \tanh \left( \frac{\pi}{2} \right) t \leq \tanh \left( \frac{\pi}{2} t \right)
\notag
\end{align}
for any $t \in [0,1]$. 
This implies inequality~\eqref{eq:ineq_equiv_concl} for $t \in [0,1]$.
For $t \geq 1$, $g(t)$ is decreasing and $t \mapsto \tanh(\frac{\pi}{2}t)$ is increasing. 
Therefore inequality~\eqref{eq:ineq_equiv_concl} also holds for  $t \geq 1$. 
Thus we obtain inequality~\eqref{eq:ineq_equiv_concl} for  any $t \geq 0$.
This completes the proof of the lemma.
\end{proof}

\end{document}
